\section{Глоссарий и перекрёстные ссылки на УГМ}
\label{app:glossary}

\paragraph{Термины.}
\begin{description}[leftmargin=0pt,itemsep=1pt]
  \item[\DM{}] регистр состояния: $\Gam\in\Dens(\mathbb{C}^7)$, $784$~байта.
  \item[Тик] одно применение $\LindOmega$ (Ур.~\ref{eq:lomega}); единственная инструкция.
  \item[Линие-шина] одна из семи линий Фано; единица интерконнекта (\S\ref{sec:fabric}).
  \item[Медиатор $k^\ast(i,j)$] третий сектор линии пары; детерминированная перемаршрутизация.
  \item[Коллапс] тождество интерконнекта, кода, инструкции и обучения как одной инцидентности (\S\ref{sec:collapse}).
  \item[Импеданс] эмуляция, которую субстрат должен добавить, чтобы запустить логическую архитектуру (\S\ref{sec:realizability}).
\end{description}

\paragraph{Использованные теоремы УГМ.} Всякое логическое заявление прослеживается
к установленному результату корпуса УГМ (\href{https://holon.sh}{holon.sh});
\TALOS{} добавляет аппаратные заявления, никакой новой теории.
\begin{center}\small
\begin{tabular}{@{}p{5.9cm}p{9.0cm}@{}}
\toprule
Результат & Использование в \TALOS{} \\
\midrule
Минимальность $7$; октонионный вывод & $N=7$ (Теор.~\ref{thm:unique}), неассоциативная ISA \\
$G_2$-ригидность (T-42a) & интерконнект $=$ орбитальная инцидентность (Теор.~\ref{thm:collapse} L4) \\
Правило отбора Фано / структурные константы $f_{ijk}$ & алгебра инструкций $=$ инцидентность линий (L3) \\
Хэмминг $[7,4,3]=$ Фано ($\Sigma$-исчисление T-224/T-225) & код $=$ топология (\S\ref{sec:ecc}) \\
Принцип третьего порядка (отпечаток Фано) & третьепорядковая маршрутизация; тай-брейк уникальности \\
$\LindOmega$, выведенный из $\Omega$ & единственная инструкция (\S\ref{sec:substrate}) \\
Метриплектическая декомпозиция (T-262) & обратимое ядро --- свободная от Ландауэра изометрия (\S\ref{sec:energy}) \\
Субстратная замкнутость (T-153/T-153a) & мультипликативность субстрата; лестница реализуемости \\
Замыкание Тегмарка (T-267), необходимость комплексной $\Gam$ (T-132) & классическое, тёплое, $784$-байтовое состояние; без охлаждения \\
Вычислительная граница (BQP) & нет ускорения по классу сложности (честная граница) \\
$\Phi$-контракция $\times\tfrac19$; время цикла $\bar\gamma\tau=\tfrac32\ln3$ & закон интеграции $9\times$; принципиальная тактовая \\
$\mathrm{SAD}_{\max}=3$ (T-142) & потолок композиции (\S\ref{sec:composition}) \\
Константы конституции ($\Rth,r_{\mathrm{stab}},\times\tfrac19$) & аппаратные лимиты (\S\ref{sec:abi}) \\
\bottomrule
\end{tabular}
\end{center}
